\documentclass[11pt,a4paper]{article}

\usepackage[utf8]{inputenc}
\usepackage[T1]{fontenc}
\usepackage[spanish]{babel}
\usepackage{amsmath,amssymb,mathtools}
\usepackage{geometry}
\usepackage{lmodern}
\usepackage{microtype}

\geometry{margin=2.4cm}
\setlength{\parindent}{0pt}
\setlength{\parskip}{0.45em}

\newcommand{\R}{\mathbb{R}}
\newcommand{\sen}{\operatorname{sen}}
\newcommand{\tr}{\operatorname{tr}}
\newcommand{\Span}{\operatorname{L}}

\begin{document}

\begin{center}
{\LARGE\bfseries Geometría Lineal}\\[0.4em]
{\Large Tema 3: Isometrías vectoriales}\\[0.4em]
{\large Soluciones}
\end{center}

\section{Ejercicio 1}
\subsection{Solución}

Sean:

\[
u=(1,2,1),
\qquad
v=(0,-1,1).
\]

Un vector ortogonal a ambos es:

\[
u\times v
=
\begin{vmatrix}
i&j&k\\
1&2&1\\
0&-1&1
\end{vmatrix}
=
(3,-1,-1).
\]

Su norma es:

\[
\|(3,-1,-1)\|
=
\sqrt{11}.
\]

Por tanto, un vector unitario buscado es:

\[
\boxed{
\frac1{\sqrt{11}}(3,-1,-1)
}
\]

\section{Ejercicio 2}
\subsection{Solución}

Para:

\[
S_1=\Span\{(1,2,3)\},
\]

un vector:

\[
(x,y,z)\in S_1^\perp
\]

debe verificar:

\[
x+2y+3z=0.
\]

Tomando:

\[
y=1,\ z=0
\]

y:

\[
y=0,\ z=1,
\]

se obtiene:

\[
\boxed{
S_1^\perp
=
\Span\{(-2,1,0),(-3,0,1)\}
}
\]

Para:

\[
S_2
=
\Span\{(0,0,1),(1,1,1)\},
\]

debe cumplirse:

\[
z=0,
\]

\[
x+y+z=0.
\]

Así:

\[
x+y=0.
\]

Por tanto:

\[
\boxed{
S_2^\perp
=
\Span\{(1,-1,0)\}
}
\]

\section{Ejercicio 3}
\subsection{Solución}

El subespacio es:

\[
S=
\{(x,y,z)\in\R^3:
2x+y+2z=0\}.
\]

Un vector normal es:

\[
n=(2,1,2).
\]

Por tanto:

\[
S^\perp
=
\Span\{(2,1,2)\}.
\]

Como:

\[
\|n\|=3,
\]

una base ortonormal es:

\[
\boxed{
\left\{
\frac13(2,1,2)
\right\}
}
\]

\section{Ejercicio 4}
\subsection{Solución}

Tomamos:

\[
v_1=(1,2,1),
\qquad
v_2=(1,1,3).
\]

El primer vector es:

\[
w_1=v_1.
\]

El segundo es:

\[
w_2
=
v_2
-
\frac{v_2\cdot w_1}{w_1\cdot w_1}w_1.
\]

Como:

\[
v_2\cdot w_1=6,
\qquad
w_1\cdot w_1=6,
\]

se obtiene:

\[
w_2
=
(1,1,3)-(1,2,1)
=
(0,-1,2).
\]

Normalizando:

\[
u_1
=
\frac1{\sqrt6}(1,2,1),
\]

\[
u_2
=
\frac1{\sqrt5}(0,-1,2).
\]

Por tanto:

\[
\boxed{
\left\{
\frac1{\sqrt6}(1,2,1),
\frac1{\sqrt5}(0,-1,2)
\right\}
}
\]

\section{Ejercicio 5}
\subsection{Solución}

Toda isometría conserva la norma.

En \(\R^3\):

\[
\|(1,1,0)\|
=
\sqrt2,
\]

\[
\|(1,1,1)\|
=
\sqrt3.
\]

Como las normas son distintas:

\[
\boxed{
\text{No existe tal isometría en }\R^3.
}
\]

En \(\R^2\):

\[
\left\|
\left(
\frac1{\sqrt2},
\frac1{\sqrt2}
\right)
\right\|
=
1
=
\|(1,0)\|.
\]

Sin embargo:

\[
(1,0)\cdot(0,1)=0
\]

mientras que:

\[
\left(
\frac1{\sqrt2},
\frac1{\sqrt2}
\right)
\cdot
(1,0)
=
\frac1{\sqrt2}\neq0.
\]

Por tanto:

\[
\boxed{
\text{Tampoco existe tal isometría en }\R^2.
}
\]

\section{Ejercicio 6}
\subsection{Solución}

Tomamos la base:

\[
B=
\{(1,1),(1,2)\}.
\]

Las imágenes son:

\[
f(1,1)=(-1,1),
\]

\[
f(1,2)=(-1,2).
\]

La matriz de cambio a la base canónica es:

\[
P=
\begin{pmatrix}
1&1\\
1&2
\end{pmatrix}.
\]

La matriz de \(f\) en la base canónica es:

\[
A
=
\begin{pmatrix}
-1&-1\\
1&2
\end{pmatrix}
P^{-1}
=
\begin{pmatrix}
-1&0\\
0&1
\end{pmatrix}.
\]

Entonces:

\[
A^tA
=
I.
\]

Por tanto:

\[
\boxed{
f\text{ es una isometría}
}
\]

Geométricamente:

\[
\boxed{
f(x,y)=(-x,y)
}
\]

es la simetría respecto del eje \(Y\).

\section{Ejercicio 7}
\subsection{Solución}

La base:

\[
B'=\{v_1,v_2\}
\]

viene dada por:

\[
v_1=e_1,
\qquad
v_2=e_2-e_1.
\]

Las imágenes son:

\[
f(v_1)=v_1+v_2=e_2,
\]

\[
f(v_2)=-2v_1-v_2=-e_1-e_2.
\]

Como:

\[
e_2=v_1+v_2,
\]

se tiene:

\[
f(e_1)=e_2.
\]

Además:

\[
e_2=v_1+v_2
\]

y, por linealidad:

\[
f(e_2)
=
f(v_1)+f(v_2)
=
-e_1.
\]

Por tanto, en la base ortonormal \(B\):

\[
M_B(f)
=
\begin{pmatrix}
0&-1\\
1&0
\end{pmatrix}.
\]

Como:

\[
M_B(f)^tM_B(f)=I,
\]

se concluye:

\[
\boxed{
f\text{ es una isometría}
}
\]

y corresponde a un giro de:

\[
\boxed{
90^\circ
}
\]

en sentido antihorario.

\section{Ejercicio 8}
\subsection{Solución}

La matriz es:

\[
A=
\begin{pmatrix}
-\frac35&-\frac45\\[1mm]
-\frac45&\frac35
\end{pmatrix}.
\]

Se verifica:

\[
A^tA=I
\]

y:

\[
\det A=-1.
\]

Por tanto, es una simetría respecto de una recta.

Buscamos los vectores fijos:

\[
(A-I)X=0.
\]

Se obtiene:

\[
2x+y=0.
\]

Por tanto:

\[
V_1
=
\Span\{(1,-2)\}.
\]

Así:

\[
\boxed{
\text{Es la simetría ortogonal respecto de la recta }y=-2x.
}
\]

\section{Ejercicio 9}
\subsection{Solución}

La recta:

\[
3x-2y=0
\]

tiene como vector director:

\[
v=(2,3).
\]

Normalizamos:

\[
u=
\frac1{\sqrt{13}}(2,3).
\]

La matriz de la simetría respecto de la recta generada por \(u\) es:

\[
A
=
2uu^t-I.
\]

Por tanto:

\[
A
=
\frac1{13}
\begin{pmatrix}
8&12\\
12&18
\end{pmatrix}
-
\begin{pmatrix}
1&0\\
0&1
\end{pmatrix}.
\]

Así:

\[
\boxed{
A=
\frac1{13}
\begin{pmatrix}
-5&12\\
12&5
\end{pmatrix}
}
\]

\section{Ejercicio 10}
\subsection{Solución}

La matriz es:

\[
A=
\begin{pmatrix}
\frac13&-\frac{2\sqrt2}{3}\\[1mm]
\frac{2\sqrt2}{3}&\frac13
\end{pmatrix}.
\]

Se verifica:

\[
\det A=1.
\]

Además tiene la forma:

\[
A=
\begin{pmatrix}
\cos\theta&-\sen\theta\\
\sen\theta&\cos\theta
\end{pmatrix}.
\]

Por tanto:

\[
\cos\theta=\frac13,
\]

\[
\sen\theta=\frac{2\sqrt2}{3}.
\]

Así:

\[
\boxed{
\theta
=
\arccos\left(\frac13\right)
}
\]

y:

\[
\boxed{
\text{La transformación es un giro antihorario de ángulo }
\arccos\left(\frac13\right).
}
\]

\section{Ejercicio 11}
\subsection{Solución}

Un giro de \(60^\circ\) en sentido horario equivale a:

\[
\theta=-60^\circ.
\]

La matriz es:

\[
A=
\begin{pmatrix}
\cos\theta&-\sen\theta\\
\sen\theta&\cos\theta
\end{pmatrix}.
\]

Por tanto:

\[
\boxed{
A=
\begin{pmatrix}
\frac12&\frac{\sqrt3}{2}\\[1mm]
-\frac{\sqrt3}{2}&\frac12
\end{pmatrix}
}
\]

\section{Ejercicio 12}
\subsection{Solución}

La matriz asociada es:

\[
A=
\begin{pmatrix}
\frac35&\frac45\\[1mm]
\frac45&-\frac35
\end{pmatrix}.
\]

Se verifica:

\[
A^tA=I.
\]

Por tanto:

\[
\boxed{
f\text{ es una transformación ortogonal}
}
\]

Además:

\[
\det A=-1,
\]

por lo que es una simetría respecto de una recta.

Los vectores fijos satisfacen:

\[
(A-I)X=0.
\]

Se obtiene:

\[
x=2y.
\]

Por tanto:

\[
\boxed{
V_1=\Span\{(2,1)\}
}
\]

y:

\[
\boxed{
f\text{ es la simetría ortogonal respecto de la recta }x=2y.
}
\]

\section{Ejercicio 13}
\subsection{Solución}

La matriz es:

\[
A=
\begin{pmatrix}
a&\frac35\\[1mm]
c&b
\end{pmatrix}.
\]

Para que sea ortogonal, sus columnas deben ser ortonormales:

\[
a^2+c^2=1,
\]

\[
\frac9{25}+b^2=1,
\]

\[
\frac35a+bc=0.
\]

De la segunda ecuación:

\[
b=\pm\frac45.
\]

Si:

\[
b=\frac45,
\]

se obtiene:

\[
(a,c)
=
\left(
\frac45,-\frac35
\right)
\]

o:

\[
(a,c)
=
\left(
-\frac45,\frac35
\right).
\]

Si:

\[
b=-\frac45,
\]

se obtiene:

\[
(a,c)
=
\left(
\frac45,\frac35
\right)
\]

o:

\[
(a,c)
=
\left(
-\frac45,-\frac35
\right).
\]

Por tanto, las cuatro posibilidades son:

\[
\boxed{
\begin{aligned}
&(a,b,c)=
\left(
\frac45,\frac45,-\frac35
\right),\\
&(a,b,c)=
\left(
-\frac45,\frac45,\frac35
\right),\\
&(a,b,c)=
\left(
\frac45,-\frac45,\frac35
\right),\\
&(a,b,c)=
\left(
-\frac45,-\frac45,-\frac35
\right).
\end{aligned}
}
\]

Las matrices con determinante \(1\) son giros y las matrices con determinante \(-1\) son simetrías respecto de una recta.

\section{Ejercicio 14}
\subsection{Solución}

Sean:

\[
R_\alpha,
\qquad
R_\beta
\]

dos rotaciones.

Entonces:

\[
R_\alpha R_\beta
=
R_{\alpha+\beta}.
\]

Por tanto:

\[
\boxed{
\text{La composición de dos rotaciones es otra rotación de ángulo }\alpha+\beta.
}
\]

Sean ahora dos simetrías respecto de rectas que forman un ángulo \(\alpha\).

La composición es un giro de ángulo:

\[
2\alpha.
\]

Por tanto:

\[
\boxed{
\text{La composición de dos simetrías es una rotación.}
}
\]

Por último:

\[
\det(R)=1,
\qquad
\det(S)=-1.
\]

Así:

\[
\det(RS)=-1.
\]

Por tanto:

\[
\boxed{
\text{La composición de una rotación y una simetría es una simetría.}
}
\]

\section{Ejercicio 15}
\subsection{Solución}

La matriz es:

\[
A=
\begin{pmatrix}
-\frac35&\frac45&0\\[1mm]
\frac45&\frac35&0\\
0&0&1
\end{pmatrix}.
\]

Se verifica:

\[
A^tA=I.
\]

Por tanto:

\[
\boxed{
A\text{ define una isometría vectorial}
}
\]

Además:

\[
\det A=-1.
\]

Los autovalores son:

\[
1,\quad1,\quad-1.
\]

El subespacio fijo viene dado por:

\[
V_1
=
\Span\{(1,2,0),(0,0,1)\}.
\]

Este plano tiene vector normal:

\[
(2,-1,0).
\]

Por tanto, su ecuación es:

\[
2x-y=0.
\]

Así:

\[
\boxed{
\text{Es la simetría ortogonal respecto del plano }2x-y=0.
}
\]

\section{Ejercicio 16}
\subsection{Solución}

El plano es:

\[
x+y-z=0.
\]

Un vector normal es:

\[
n=(1,1,-1).
\]

La matriz de la simetría respecto del plano es:

\[
A
=
I
-
2\frac{nn^t}{n^tn}.
\]

Como:

\[
n^tn=3,
\]

se obtiene:

\[
\boxed{
A
=
\frac13
\begin{pmatrix}
1&-2&2\\
-2&1&2\\
2&2&1
\end{pmatrix}
}
\]

Toda rotación en \(\R^3\) tiene determinante \(1\).

Por tanto, la composición de dos giros tiene determinante:

\[
1.
\]

Sin embargo:

\[
\det A=-1.
\]

Así:

\[
\boxed{
\text{No es posible obtener esta simetría como composición de dos giros.}
}
\]

\section{Ejercicio 17}
\subsection{Solución}

La matriz es:

\[
A=
\begin{pmatrix}
0&1&0\\
1&0&0\\
0&0&-1
\end{pmatrix}.
\]

Se verifica:

\[
A^tA=I,
\qquad
\det A=1.
\]

Los autovalores son:

\[
1,-1,-1.
\]

El subespacio de vectores fijos es:

\[
V_1
=
\Span\{(1,1,0)\}.
\]

Por tanto:

\[
\boxed{
\text{Es un giro de }180^\circ
\text{ alrededor de la recta }
\Span\{(1,1,0)\}.
}
\]

\section{Ejercicio 18}
\subsection{Solución}

Sea:

\[
u=
\frac1{\sqrt3}(1,1,1).
\]

La simetría respecto de la recta generada por \(u\) deja fija la componente paralela a \(u\) y cambia de signo la componente perpendicular.

Por tanto:

\[
A
=
2uu^t-I.
\]

Así:

\[
\boxed{
A
=
\frac13
\begin{pmatrix}
-1&2&2\\
2&-1&2\\
2&2&-1
\end{pmatrix}
}
\]

\section{Ejercicio 19}
\subsection{Solución}

La matriz es:

\[
A=
\begin{pmatrix}
0&0&1\\
1&0&0\\
0&1&0
\end{pmatrix}.
\]

Se verifica:

\[
A^tA=I,
\qquad
\det A=1.
\]

Además:

\[
A(1,1,1)=(1,1,1).
\]

Por tanto:

\[
V_1
=
\Span\{(1,1,1)\}.
\]

Como:

\[
\tr A=0,
\]

para un giro en \(\R^3\):

\[
1+2\cos\theta=0.
\]

Así:

\[
\cos\theta=-\frac12.
\]

Por tanto:

\[
\boxed{
\theta=120^\circ
}
\]

y:

\[
\boxed{
\text{Es un giro de }120^\circ
\text{ alrededor de la recta }
\Span\{(1,1,1)\}.
}
\]

\section{Ejercicio 20}
\subsection{Solución}

El eje tiene vector unitario:

\[
u=
\frac1{\sqrt3}(1,1,1).
\]

Aplicamos la fórmula de Rodrigues para:

\[
\theta=60^\circ.
\]

La matriz obtenida es:

\[
\boxed{
A=
\frac13
\begin{pmatrix}
2&-1&2\\
2&2&-1\\
-1&2&2
\end{pmatrix}
}
\]

Se verifica:

\[
A(1,1,1)=(1,1,1).
\]

\section{Ejercicio 21}
\subsection{Solución}

La matriz es:

\[
A=
\begin{pmatrix}
-\frac12&-\frac{\sqrt2}{2}&\frac12\\[1mm]
\frac{\sqrt2}{2}&0&\frac{\sqrt2}{2}\\[1mm]
\frac12&-\frac{\sqrt2}{2}&-\frac12
\end{pmatrix}.
\]

Se verifica:

\[
A^tA=I,
\qquad
\det A=-1.
\]

Además:

\[
\tr A=-1.
\]

Para una isometría impropia de tipo giro más simetría:

\[
\tr A
=
-1+2\cos\theta.
\]

Por tanto:

\[
-1=-1+2\cos\theta
\]

y:

\[
\boxed{
\theta=\frac{\pi}{2}
}
\]

El subespacio asociado al autovalor \(-1\) es:

\[
V_{-1}
=
\Span\{(-1,0,1)\}.
\]

Por tanto, el eje es:

\[
\boxed{
r=\Span\{(-1,0,1)\}
}
\]

y el plano de simetría es:

\[
V_{-1}^\perp.
\]

Su ecuación es:

\[
-x+z=0.
\]

Así:

\[
\boxed{
\text{Es la composición de un giro de }90^\circ
\text{ alrededor de }r
\text{ y una simetría respecto del plano }z=x.
}
\]

\section{Ejercicio 22}
\subsection{Solución}

La simetría respecto del plano:

\[
z=0
\]

tiene matriz:

\[
S=
\begin{pmatrix}
1&0&0\\
0&1&0\\
0&0&-1
\end{pmatrix}.
\]

El giro de \(120^\circ\) alrededor del eje \(OZ\) tiene matriz:

\[
R=
\begin{pmatrix}
-\frac12&-\frac{\sqrt3}{2}&0\\[1mm]
\frac{\sqrt3}{2}&-\frac12&0\\
0&0&1
\end{pmatrix}.
\]

Como ambas transformaciones conmutan:

\[
A=RS=SR.
\]

Por tanto:

\[
\boxed{
A=
\begin{pmatrix}
-\frac12&-\frac{\sqrt3}{2}&0\\[1mm]
\frac{\sqrt3}{2}&-\frac12&0\\
0&0&-1
\end{pmatrix}
}
\]

\section{Ejercicio 23}
\subsection{Solución}

Se sabe que:

\[
f(1,0,0)=(-1,0,0).
\]

Por tanto:

\[
(1,0,0)\in V_{-1}.
\]

Además:

\[
\det A=-1
\]

y:

\[
\tr A=1.
\]

Una isometría de \(\R^3\) con determinante \(-1\), traza \(1\) y un autovalor \(-1\) es una simetría respecto de un plano.

Como:

\[
V_{-1}
=
\Span\{(1,0,0)\},
\]

el plano fijo es:

\[
V_{-1}^\perp
=
\{(x,y,z):x=0\}.
\]

Por tanto:

\[
\boxed{
\text{Es la simetría ortogonal respecto del plano }x=0.
}
\]

\section{Ejercicio 24}
\subsection{Solución}

La aplicación es:

\[
f(x,y,z)
=
\frac13
(-x+2y+2z,\,
2x-y+2z,\,
2x+2y-z).
\]

Su matriz es:

\[
A=
\frac13
\begin{pmatrix}
-1&2&2\\
2&-1&2\\
2&2&-1
\end{pmatrix}.
\]

Se verifica:

\[
A^tA=I.
\]

Por tanto:

\[
\boxed{
f\text{ es una isometría}
}
\]

Además:

\[
\det A=1.
\]

Los autovalores son:

\[
1,-1,-1.
\]

El subespacio fijo es:

\[
V_1
=
\Span\{(1,1,1)\}.
\]

Por tanto:

\[
\boxed{
f\text{ es la simetría respecto de la recta }
\Span\{(1,1,1)\}
}
\]

o, equivalentemente:

\[
\boxed{
f\text{ es un giro de }180^\circ
\text{ alrededor de dicha recta.}
}
\]

\section{Ejercicio 25}
\subsection{Solución}

Queremos:

\[
e_1\longmapsto e_2,
\qquad
e_2\longmapsto e_3,
\qquad
e_3\longmapsto e_1.
\]

La matriz correspondiente es:

\[
A=
\begin{pmatrix}
0&0&1\\
1&0&0\\
0&1&0
\end{pmatrix}.
\]

Se verifica:

\[
A(1,1,1)=(1,1,1).
\]

Por tanto, el eje es:

\[
\boxed{
r=\Span\{(1,1,1)\}
}
\]

Además:

\[
\tr A=0.
\]

Así:

\[
1+2\cos\theta=0.
\]

Por tanto:

\[
\boxed{
\theta=120^\circ
}
\]

\section{Ejercicio 26}
\subsection{Solución}

La base es:

\[
B=
\{(1,1,0),(0,1,0),(0,0,1)\}.
\]

La matriz de cambio de \(B\) a la base canónica es:

\[
P=
\begin{pmatrix}
1&0&0\\
1&1&0\\
0&0&1
\end{pmatrix}.
\]

La matriz de \(f\) en \(B\) es:

\[
A_B=
\begin{pmatrix}
\frac35&0&-\frac45\\[1mm]
-\frac85&-1&\frac45\\[1mm]
\frac45&0&\frac35
\end{pmatrix}.
\]

En la base canónica:

\[
A_c
=
PA_BP^{-1}
=
\begin{pmatrix}
\frac35&0&-\frac45\\[1mm]
0&-1&0\\[1mm]
\frac45&0&\frac35
\end{pmatrix}.
\]

Se verifica:

\[
A_c^tA_c=I.
\]

Por tanto:

\[
\boxed{
f\text{ es una transformación ortogonal}
}
\]

Además:

\[
\det A_c=-1.
\]

El autovalor \(-1\) tiene subespacio:

\[
V_{-1}
=
\Span\{(0,1,0)\}.
\]

Así, el eje es el eje \(OY\) y el plano perpendicular es:

\[
y=0.
\]

La traza es:

\[
\tr A_c
=
\frac15.
\]

Para una isometría de tipo \(5\):

\[
\tr A_c
=
-1+2\cos\theta.
\]

Por tanto:

\[
\cos\theta=\frac35.
\]

Así:

\[
\boxed{
\theta
=
\arccos\left(\frac35\right)
}
\]

Por tanto:

\[
\boxed{
\text{Es la composición de un giro alrededor del eje }OY
\text{ de ángulo }\arccos\left(\frac35\right)
\text{ y una simetría respecto del plano }y=0.
}
\]

\section{Ejercicio 27}
\subsection{Solución}

El plano es:

\[
S=
\{(x,y,z)\in\R^3:x-y=0\}.
\]

Una base ortonormal de \(S\) es:

\[
u_1=
\frac1{\sqrt2}(1,1,0),
\qquad
u_2=(0,0,1).
\]

Un vector unitario de \(S^\perp\) es:

\[
n=
\frac1{\sqrt2}(1,-1,0).
\]

Sobre \(S\), un giro de ángulo:

\[
\pi
\]

actúa como:

\[
u_1\mapsto-u_1,
\qquad
u_2\mapsto-u_2.
\]

Sobre \(S^\perp\) existen dos posibilidades:

\[
n\mapsto n
\]

o:

\[
n\mapsto-n.
\]

En el primer caso se obtiene:

\[
\boxed{
A_1=
\begin{pmatrix}
0&-1&0\\
-1&0&0\\
0&0&-1
\end{pmatrix}
}
\]

que es un giro de \(180^\circ\) alrededor de la recta:

\[
S^\perp
=
\Span\{(1,-1,0)\}.
\]

En el segundo caso:

\[
\boxed{
A_2=-I
=
\begin{pmatrix}
-1&0&0\\
0&-1&0\\
0&0&-1
\end{pmatrix}
}
\]

que es la simetría respecto del origen.

Por tanto:

\[
\boxed{
A_1
\text{ y }
A_2
\text{ son todas las posibilidades.}
}
\]

\section{Ejercicio 28}
\subsection{Solución}

La matriz es:

\[
A=
\begin{pmatrix}
a&a&b\\
a&a&-b\\
-b&b&0
\end{pmatrix}.
\]

Para que sea ortogonal, sus columnas deben ser ortonormales.

De la tercera columna:

\[
2b^2=1.
\]

Por tanto:

\[
b=\pm\frac1{\sqrt2}.
\]

Además:

\[
( a,a,-b)\cdot(a,a,b)
=
2a^2-b^2
=
0.
\]

Así:

\[
2a^2=\frac12.
\]

Por tanto:

\[
a=\pm\frac12.
\]

Luego:

\[
\boxed{
a=\pm\frac12,
\qquad
b=\pm\frac1{\sqrt2}
}
\]

con signos independientes.

Si:

\[
a=\frac12,
\]

se tiene:

\[
\det A=1,
\qquad
\tr A=1.
\]

Por tanto:

\[
1+2\cos\theta=1
\]

y:

\[
\boxed{
\theta=\frac{\pi}{2}
}
\]

El eje es:

\[
\boxed{
V_1=\Span\{(1,1,0)\}.
}
\]

Así, para ambos signos de \(b\), es un giro de \(90^\circ\), cambiando el sentido según el signo de \(b\).

Si:

\[
a=-\frac12,
\]

entonces:

\[
\det A=-1,
\qquad
\tr A=-1.
\]

Así:

\[
-1+2\cos\theta=-1
\]

y:

\[
\boxed{
\theta=\frac{\pi}{2}.
}
\]

Además:

\[
V_{-1}
=
\Span\{(1,1,0)\}.
\]

Por tanto:

\[
\boxed{
\text{Es una composición de un giro de }90^\circ
\text{ alrededor de }
\Span\{(1,1,0)\}
\text{ y una simetría respecto del plano }x+y=0.
}
\]

El signo de \(b\) determina el sentido del giro.

\end{document}
